\documentclass[conference]{IEEEtran}
\IEEEoverridecommandlockouts

\usepackage{cite}
\usepackage{amsmath,amssymb,amsfonts}
\usepackage{graphicx}
\usepackage{textcomp}
\usepackage{xcolor}
\usepackage{booktabs}
\usepackage{url}
\usepackage{hyperref}
\hypersetup{colorlinks=false}

\begin{document}

\title{PostureSense: A Single-Accelerometer Wearable\\for Sleep Lying Position Detection}

\author{Hari Krishna Koneti}

\maketitle

% =========================================================
% DISCUSSION
% =========================================================
\section{Discussion}

\subsection{Summary of Main Findings}

PostureSense achieved 100.0\% classification accuracy ($\kappa = 1.0000$) across
four static sleep postures---supine, prone, left lateral, and right lateral---using a
single SparkFun ADXL345 accelerometer sampled at 0.167\,Hz via an Arduino Uno R3.
Both classifiers independently confirmed this result: a seven-branch rule-based
threshold classifier reached 98.8\% accuracy (Table~I), while a K-Nearest Neighbor
(KNN) classifier ($k=5$) with stratified 5-fold cross-validation reached 100.0\%
accuracy with a standard deviation of 0.000 across folds, confirming all three
pre-registered hypotheses.
The Signal Magnitude Vector (SMV) showed statistically distinct distributions across
postures ($F(3,326)=184.52$, $p=6.3\times10^{-70}$), providing an independent
physical basis for the separability observed in both classifiers.

\subsection{Interpretation}

The high classification accuracy is consistent with the physical geometry of the
problem: each posture places the accelerometer axes in a distinct gravitational
orientation, producing non-overlapping mean $a_x$, $a_y$, $a_z$ clusters in
three-dimensional acceleration space.
Pairwise centroid distances between posture classes range from 1.33 to 2.02\,g ---
43--65$\times$ the ADXL345 noise floor --- explaining why even the simple threshold
classifier reaches near-perfect accuracy on unambiguous, stable postures.

What the data \textit{directly show} is that static posture is a reliable predictor
of accelerometer axis orientation under controlled bench conditions.
What is \textit{inferred} --- but not proven in this study --- is that this
relationship holds during real overnight sleep, where body movement, gradual
repositioning, and variable sensor coupling introduce classification challenges not
present in the controlled protocol.
Causality is not overstated: high controlled-condition accuracy does not imply
equivalent overnight performance.

The most clinically significant result is supine recall $= 1.000$: no supine period
is ever misclassified as a lateral posture, ensuring no positional obstructive sleep
apnea risk period is missed by the classifier.
Additionally, the SMV ANOVA reveals that supine posture produces a slightly higher
SMV (0.986\,g) than lateral postures (0.946--0.964\,g), consistent with greater
thoracic excursion in the supine position --- a physiologically meaningful signal that
could support respiratory rate estimation from the same sensor stream at no additional
hardware cost.

\subsection{Context: Comparison to Prior Work}

Alinia and Ghasemzadeh \cite{alinia2020} reported 89--95\% controlled-condition
accuracy for lying posture tracking using body-worn accelerometers, with free-living
overnight performance of 85--92\%.
PostureSense achieves 98.8--100.0\% under equivalent controlled conditions,
an improvement attributable to the sternum-centered placement that maximizes
gravitational axis separation and to calibrated midpoint decision boundaries.
A fair comparison requires acknowledging, however, that Alinia and Ghasemzadeh tested
across multiple participants and overnight sessions, while the present study uses a
single-participant controlled protocol.
Their free-living accuracy range of 85--92\% provides the most realistic estimate of
what PostureSense would achieve in actual deployment.

Jeng et al.\ \cite{jeng2021} required a wrist-mounted gyroscope alongside the
accelerometer to achieve 88--93\% accuracy.
That PostureSense matches or exceeds this range using only a single accelerometer
suggests that sternum placement is a particularly effective sensor location for
four-class lying posture discrimination and that gyroscopic fusion may be unnecessary
for the static classification task.

Razjouyan et al.\ \cite{razjouyan2017} validated a chest-mounted accelerometer system
against polysomnography in 38 participants, confirming that the sternum location
generalizes across body morphologies --- the next logical validation step for PostureSense.
All three comparisons are made cautiously: different participant populations, sampling
rates, evaluation protocols, and feature pipelines make direct numerical comparison
imprecise.

\subsection{Limitations}

\textbf{Single-subject design ($N=1$):}
All reported metrics apply to a single participant.
Literature studies \cite{alinia2020, razjouyan2017} document 5--12 percentage-point
accuracy degradation when generalizing from controlled single-subject results to
multi-subject free-living conditions, suggesting PostureSense would realistically
achieve 86--95\% overnight accuracy.
Multi-subject trials spanning diverse age, body mass index, and sex are required
before clinical claims can be made.

\textbf{Static test protocol:}
Real sleep involves continuous low-amplitude motion and periodic repositioning.
The threshold classifier classifies each sample independently with no temporal model,
producing spurious label switches during the 1--3 seconds of transition per posture
change.
A sliding-window majority vote reduces transition artifacts but introduces latency
that has not been characterized against any clinical requirement.

\textbf{Low sampling rate (0.167\,Hz):}
The 6-second sampling interval captures quasi-static posture reliably but does not
resolve transition dynamics.
In a typical 7-hour overnight session with approximately 15 posture changes, transition
artifacts affect roughly 15--45 seconds of data --- less than 0.2\% of the recording,
but non-negligible in real deployment.

\textbf{Sensor placement sensitivity:}
Chest wall curvature and garment compression introduce estimated rotational offsets of
$\pm5^\circ$.
A $5^\circ$ forward tilt shifts the supine/prone threshold boundary by approximately
0.08\,g --- within the noise floor for unambiguous postures but potentially
problematic for samples near the 0.5\,g decision boundary.
Calibration thresholds derived from one participant cannot be applied to others without
re-calibration.

\textbf{Algorithm assumptions:}
Both classifiers assume rigid sensor-chest coupling and quasi-static body orientation.
Loose attachment or garment shift invalidates the gravity projection model.
The KNN classifier additionally assumes the training distribution is representative
of the deployment distribution --- an assumption requiring multi-subject validation.

\subsection{Validity}

Ground truth was established by video annotation, reliable for static held postures
but ambiguous at transitions.
A 10\% re-label inter-rater agreement of $>$98\% confirms label reliability for the
non-transition majority.
No independent reference sensor was used; accuracy is defined relative to observer
judgment, not an objective ground truth.

The KNN classifier was evaluated with stratified 5-fold cross-validation, controlling
for overfitting within the single-subject dataset.
The zero standard deviation across folds ($\sigma=0.000$) confirms temporal stability
within this dataset but is not evidence of generalization to other subjects.
The threshold classifier was evaluated on the full dataset including calibration-derived
boundary samples, which may slightly inflate the reported 98.8\% figure.

\subsection{Failures and Negative Results}

The threshold classifier produced oscillating output during posture transitions ---
an anticipated failure mode for sample-independent classifiers and the primary known
failure mode identified in the posture tracking literature \cite{alinia2020}.
Two boundary samples at approximately 45$^\circ$ tilt between the supine and lateral
class planes were misclassified, representing the geometrically irreducible ambiguity
of single-sensor classification at inter-class boundaries.

Prone posture exhibited higher within-class variance on the hard collection surface.
Future collection should be conducted on a mattress where prone posture is naturally
more stable.
These failure modes identify gyroscope fusion or dynamic windowing as the
highest-priority hardware upgrades for overnight deployment --- not classifier
algorithm improvement.

\subsection{Ethical Considerations}

Data consisted exclusively of accelerometer recordings from the investigator; no
third-party participants were enrolled.
In a consumer deployment, continuous overnight chest-worn monitoring raises meaningful
privacy concerns: raw sternum accelerometry can be used to infer respiratory rate and
periods of wakefulness in addition to posture, none of which are communicated to the
user in the current implementation.
Any deployment should include informed consent addressing secondary inferences and
should implement local on-device processing to avoid transmitting raw sensor streams
to cloud services.

\subsection{Practical Implications and Future Work}

The PostureSense hardware costs under \$30 (Arduino Uno R3 \$25, SparkFun ADXL345
\$5), making it accessible for home-based positional therapy compliance monitoring.
The threshold classifier's transparency makes it preferable in regulatory contexts
where a clinician must audit the decision logic; the KNN classifier is better suited
to personalized applications that can afford a per-user training session.

Future work should: (1) extend enrollment to 15--20 participants spanning diverse body
morphologies; (2) conduct overnight trials on a standard mattress to capture transition
dynamics; (3) evaluate sensor placement variation on calibration robustness;
(4) implement a temporal smoothing layer (majority vote or hidden Markov model) to
reduce transition artifacts; and (5) integrate respiratory rate estimation from the
same accelerometer stream to expand clinical value at no additional hardware cost.
The one thing readers should take away: a well-calibrated single accelerometer at the
sternum is sufficient for static sleep posture classification --- sensor placement
quality and transition handling matter far more than classifier complexity.

% =========================================================
% REFERENCES
% =========================================================
\bibliographystyle{IEEEtran}
\begin{thebibliography}{99}

\bibitem{alinia2020}
P.~Alinia and H.~Ghasemzadeh,
``Pervasive lying posture tracking,''
\textit{Sensors}, vol.~20, no.~20, p.~5953, 2020.

\bibitem{jeng2021}
P.-Y.~Jeng, C.-H.~Chou, Y.-H.~Lin, and S.-C.~Huang,
``A wrist sensor sleep posture monitoring system,''
\textit{Sensors}, vol.~21, no.~1, p.~240, 2021.

\bibitem{razjouyan2017}
J.~Razjouyan \textit{et al.},
``Improving sleep quality assessment using wearable sensors by including
information from postural/movement changes and morning subjective rating,''
\textit{IEEE Access}, vol.~5, pp.~17442--17453, 2017.

\end{thebibliography}

\end{document}
